\section{Evaluation}\label{sec:eval}
Our evaluation is designed to answer three key research questions. 

\begin{enumerate}[label=(\arabic*)]
    \item Can SHMuzz discover previously unknown double-fetch vulnerabilities in commercial Trusted Applications?
    \item Can the vulnerabilities identified by SHMuzz be reproduced on consumer off-the-shelf devices?
    \item Are Trusted Applications written in memory-safe languages also susceptible to double-fetch attacks?
\end{enumerate}

\subsection{Fuzzing Commercial TAs}



\begin{table*}[ht]
    \scriptsize
    \centering
    %\begin{adjustbox}{max width=\columnwidth}
    \begin{tabular}{lrrrrrrr}
        \multicolumn{3}{c}{Dataset} & \multicolumn{3}{c}{Exploration} & \makecell[c]{Double Fetch\\Snapshot Fuzzing} & \makecell[c]{Distillation}\\
        \cmidrule(lr){1-3} \cmidrule(lr){4-6} \cmidrule(lr){7-7} \cmidrule(lr){8-8} 
        \makecell[c]{TEE}&
        \makecell[c]{\# TAs}&
        \makecell[c]{\# TAs Harnessed}& 
        \makecell[c]{\# TAs with\\Overl. Fetches}& 
        \makecell[c]{\# Overl. Fetches}&
        \makecell[c]{\# Overl. Fetches merged} &
        \makecell[c]{\# Crashes}&
        \makecell[c]{\# Crashes distilled}\\
        \midrule
        TEEGris    & 33   & 9   & 7  & 939         & 35      & 74 & 37\\
        QSEE       & 4    & 3   & 2  & 1,232,565   & 2,345   & 12 & 6\\
        MiTEE      & 12   & 10  & 8  & 35,408      & 1,581   & 64 & 15\\
        Beanpod    & 11   & 5   & 4  & 18,785      & 2,576   & 18 & 9\\
        \midrule
        \textbf{all}& 60          & 27          & 21         &  1,287,697          & 6,537          & 168 & 67\\
        \bottomrule
    \end{tabular}
    %\end{adjustbox}
    \caption{
        \yiwen{TODO}
    }
    \label{tab:dataset}
\end{table*}


\subsection{Reproducing DF Vulnerabilities on Phones}
Our implementation of SHMuzz assumes that {\tt memref} buffers are backed by shared memory. However, as discussed in Section~\ref{sec:bak}, some TEE implementations may copy the buffer instead. To confirm that the issues identified by SHMuzz are real security vulnerabilities, We reproduce the discovered double-fetch attacks on consumer off-the-shelf (COTS) devices (see Table~\ref{}). All test devices are rooted and updated to the latest firmware.

I develop a proof-of-concept (PoC) normal-world client application to trigger the vulnerabilities on real devices. The PoC performs three key operations: obtaining a pointer to shared memory, invoking the vulnerable TA functionality, and racing the shared-memory access.

\begin{table*}[t]
    \small
    \centering
    %\begin{adjustbox}{max width=\columnwidth}
    \begin{tabular}{llllc}
        TEE&
        TA UUID &
        TA Name & 
        DF Vulnerability &
        \makecell{Reproduced\\on Device} \\
        \midrule
        MiTEE      & 377ee4e8-af0e-474f-a9d636a9268fe85c   & SoterApp & stack overflow    & \checkmark \\
        TEEGris    & 00000000-0000-0000-0000-4662436b6d52  & FbSkmR & oob read          & \checkmark \\
        MiTEE,QSEE& 3d08821c-33a6-11e6-a1fa089e01c83aa2   & VSIMApp & oob write         & \checkmark* \\
        MiTEE,QSEE& 3d08821c-33a6-11e6-a1fa089e01c83aa2   & VSIMApp & double free       & \checkmark* \\
        MiTEE      & 88ce8e6b-8646-4092-bb78faf5b55ff4df   & Mlipay & oob read          & \checkmark \\
        Beanpod   & 08010203000000000000000000000000       & ifaa-key & oob read          & \xmark \dag \\
        \bottomrule
    \end{tabular}
    %\end{adjustbox}
    \caption{
    The vulnerabilities in commercial TAs found with SHMuzz.
    (*) Only on MiTEE, because we did not have a QSEE testing device supporting this TA.
    (\dag) None of our testing devices supports the TA. 
    }
    \label{tab:vulns}
\end{table*}

The PoC communicates with TAs using the vendor-provided GlobalPlatform client library, which interfaces with the TEE-specific kernel driver. This allows most of the PoC code to be reused across different TEEs. To allocate shared memory, the PoC calls {\tt TEEC\_AllocateSharedMemory}. However, on MiTEE and TEEGris, this API returns a pointer to an intermediate buffer whose contents are copied into the actual shared memory by the client library. To access the true shared-memory region, We reverse engineer the client libraries and directly extract the internal pointer. Listing~\ref{} (lines 13$\sim$19) illustrates this process for TEEGris, where the PoC retrieves a TEE-specific internal structure and uses it to locate the real shared-memory address.

After obtaining the shared-memory pointer, the PoC spawns a separate thread that continuously modifies the buffer while the TA is executing. This asynchronous update races the TA’s accesses and eventually triggers the double-fetch condition. As shown in Listing~\ref{} (lines 1$\sim$6), the racing thread alternates between two values: mod1, which ensures the TA’s initial check succeeds, and mod2, which corrupts the data to cause a failure. If mod2 is written between the two fetches, the TA observes inconsistent values, thereby triggering a TOCTTOU vulnerability.

The PoC then invokes the vulnerable TA function with carefully constructed arguments. These inputs are derived from the fuzzer-generated seed that originally exposed the overlapping fetch (Listing~\ref{}, lines 24$\sim$25). We repeatedly execute the PoC until a crash is observed, which is detected by checking the return code of {\tt TEEC\_InvokeCommand}. Using this approach, We successfully reproduced five of the six discovered double-fetch vulnerabilities on physical devices.


\begin{figure}[h]
    \definecolor{pgreen}{rgb}{0,0.5,0}
\definecolor{pblue}{rgb}{0.13,0.13,1}
\definecolor{pred}{rgb}{0.9,0,0}
\definecolor{pgrey}{rgb}{0.46,0.45,0.48}
\definecolor{light-gray}{gray}{0.975}
\definecolor{darkgreen}{rgb}{0,0.5,0}
\definecolor{vscodepurple}{RGB}{128, 0, 128}
\definecolor{codepurple}{rgb}{0.58,0,0.82}

\lstset{escapechar=@,
                backgroundcolor = \color{light-gray},
                linebackgroundcolor={%
                \ifnum\value{lstnumber}>0
                    \ifnum\value{lstnumber}<30
                        \color{light-gray}
                    \fi
                \fi
                %\ifnum\value{lstnumber}>10
                %    \ifnum\value{lstnumber}<12
                %        \color{teal!40}
                %    \fi
                %\fi
                \ifnum\value{lstnumber}>3
                    \ifnum\value{lstnumber}<5
                        \color{blue!20}
                    \fi
                \fi 
                \ifnum\value{lstnumber}>6
                    \ifnum\value{lstnumber}<8
                        \color{red!20}
                    \fi
                \fi
                },
                tabsize=2,
                basicstyle=\fontencoding{T1}\selectfont\footnotesize,
                %basicstyle=\fontfamily{lmvtt}\selectfont\footnotesize,
                commentstyle=\color{darkgreen!80},
                keywordstyle=\color{pblue},
                stringstyle=\color{pred},
                %stringstyle=\color{vscodepurple},
                frame = single,
                showstringspaces=false,
                language=c,
                xleftmargin=5pt, 
                xrightmargin=5pt, 
                %framexleftmargin=-5pt,
                numbers=left,                  % Line numbers
                numberstyle=\tiny\color{black}, % Line number style
                numbersep=4pt, %moredelim=[is][\textbf]{@}{@}
                %escapeinside={(*@}{@*)}
}
\center

\begin{lstlisting}[columns=fullflexible]
void combine_file_name(char* part1, char* shm){
  char out_path[0x80];
  size_t p1_len = strlen(part1);
  size_t p2_len = strlen(shm); 
  if(p1_len + p2_len > 0x80) return;
  strcpy(out_path, p1);
  strcpy(out_path + p1_len, shm);
  ...
}
\end{lstlisting}

    \captionof{lstlisting}{
        The simplified code of the double fetch stack overflow vulnerability in the 
        \texttt{377ee4e8-af0e-474f-a9d636a9268fe85c} TA. 
        The \texttt{shm} argument of the function points to the shared \texttt{memref} buffer. 
    }
    \label{lst:soterbug}
\end{figure}

{\bf Case Study: Stack Overflow in MiTEE}.
The affected Trusted Application contains a double-fetch vulnerability that results in a stack-based buffer overflow. As shown in Listing~\ref{}, the TA first retrieves the length of a null-terminated string stored in shared memory. During the second fetch, it copies the same string into a fixed-size stack buffer.
To exploit this flaw, the PoC’s racing thread initializes the shared memory with a string longer than 0x80 bytes. It then repeatedly toggles the first byte of the buffer between zero and its original value. If the race succeeds, {\tt strlen} is computed on an empty string, while {\tt strcpy} later copies the full, oversized string. This mismatch causes the copy operation to overflow the stack buffer, triggering memory corruption.



\begin{table}[t]
    \small
    \centering
    %\begin{adjustbox}{max width=\columnwidth}
    \begin{tabular}{llc}
        TEE&
        Phone &
        Zero-Copy Shm\\
        \midrule
        TEEGris    & Samsung A56 5G   & \checkmark \\
        TEEGris    & Samsung A16   & \checkmark \\
        TEEGris    & Samsung S22   & \checkmark \\
        TEEGris    & Samsung S10   & \checkmark \\
        QSEE       & OnePlus 12R   & \checkmark \\
        QSEE       & OnePlus Nord 5  & \checkmark \\
        MiTEE      & Xiaomi 15T   & \checkmark \\
        MiTEE      & Xiaomi Redmi Note 13 5G   & \checkmark \\
        Beanpod    & Xiaomi Redmi Note 11s 5G  & \xmark \\
        Beanpod    & Xiaomi Redmi 9A  & \xmark \\
        \bottomrule
    \end{tabular}
    %\end{adjustbox}
    \caption{
        The devices we used to reproduce double fetches. All devices were updated to the newest firmware version. 
    }
    \label{tab:phones}
\end{table}

{\bf Case Study: Out-of-Bounds Read in TEEGris}.
The affected TEEGris Trusted Application contains a double-fetch vulnerability that leads to an out-of-bounds read on a shared-memory input buffer. The TA expects the buffer to store serialized length–value data. As shown in Listing~\ref{}, it first verifies that the declared length does not exceed the buffer size. It then fetches the length again and uses it to copy data from shared memory into a local buffer.
To exploit this flaw, the PoC’s racing thread repeatedly alternates the first {\tt size\_t} value in shared memory between zero and a very large value. If the modification occurs between the two fetches, the second value causes {\tt memcpy} to read beyond the bounds of the buffer, triggering an out-of-bounds memory access.

{\bf Double Fetches in Other TEEs}.
For TEEs where on-device crashes could not be reproduced, namely Kinibi, Beanpod, and QSEE, either because SHMuzz did not identify a crashing input or because the affected TAs were unavailable on the test devices, an additional analysis was conducted to determine whether zero-copy shared memory is supported.
The overlapping fetches identified by SHMuzz were manually examined, and those with observable side effects under shared-memory races were selected. Using this method, zero-copy shared memory was confirmed for both Kinibi and QSEE (see Appendix~\ref{}). In contrast, Beanpod does not appear to expose shared memory between the normal and secure worlds.
To further validate this observation, a previously published exploit was used to obtain code execution in a Beanpod TA. The injected shellcode repeatedly read from a {\tt memref} buffer and logged its value. Although the normal-world client modified the buffer concurrently, no changes were observed by the TA, indicating that Beanpod relies on memory copying rather than direct sharing. A summary of zero-copy support across TEEs is provided in Table~\ref{}.